\chapter{Anforderungen und Randbedingungen}
\label{ch:anforderungen}

Die Anforderungen an Modul~4 wurden zu Beginn in einem eigenen Anforderungsdokument
festgehalten und während der Arbeit fortgeschrieben
\quelle{Requierments/requierments.md}. Dieses Kapitel fasst sie zusammen und ordnet sie den
Hürden aus \cref{sec:huerden} zu.

% ==============================================================================
\section{Funktionale Anforderungen}
\label{sec:anf-funktional}

\begin{description}
  \item[F1 Aufgabe] Die Policy greift das Bauteil aus beliebiger Lage im Arbeitsbereich und
    fährt es an den Übergabepunkt über der Prüfvorrichtung. Dort übergibt sie an Modul~2.
  \item[F2 Beobachtung] Die Policy arbeitet mit den Bildern einer Handgelenkkamera und der
    Draufsichtkamera sowie dem Roboterzustand. Tiefendaten stehen nicht zur Verfügung.
  \item[F3 Aufnahme] Demonstrationen werden mit möglichst wenig Handarbeit aufgenommen.
    Insbesondere soll eine einmal hergestellte Lage mehrfach und neue Lagen ohne erneutes
    Teachen aufgenommen werden können.
  \item[F4 Training] Das Training läuft lokal auf einem Rechner der Gruppe.
  \item[F5 Ausführung] Die Policy läuft in Echtzeit mit festem Takt am realen Roboter, in der
    \ac{vm} der Steuerung und in einer eigenen Simulation.
  \item[F6 Einbindung] Modul~4 lässt sich aus dem Ablauf und der Bedienoberfläche des
    Gesamtsystems starten.
\end{description}

% ==============================================================================
\section{Nichtfunktionale Anforderungen}
\label{sec:anf-nichtfunktional}

\begin{description}
  \item[N1 Lokal und autark] Weder Daten noch Modelle werden in einen Cloud-Dienst
    hochgeladen. Alle Werkzeuge laufen auf eigenen Rechnern.
  \item[N2 Hardwarefrei entwickelbar] Zu Beginn der Arbeit war weder der Roboter noch eine
    Kamera verfügbar. Alles, was nicht zwingend an der Hardware hängt, muss ohne sie lauffähig
    und getestet sein (\cref{ch:architektur}).
  \item[N3 Austauschbar] Roboter und Kameras sollen ersetzt werden können, ohne Aufnahme,
    Training oder Ausführung zu ändern.
  \item[N4 Sicher] Jeder Befehl der gelernten Policy wird vor dem Senden geprüft. Der
    Hardware-Not-Halt bleibt die letzte Instanz (\cref{ch:sicherheit}).
  \item[N5 Nachvollziehbar] Jede Aufnahme, jedes Modell und jeder Versuch ist mit seinen
    Einstellungen protokolliert, sodass Ergebnisse einem Datenstand zugeordnet werden können.
\end{description}

% ==============================================================================
\section{Randbedingungen der Anlage}
\label{sec:randbedingungen}

Einige Eigenschaften der Anlage schränken die Lösung ein und haben Designentscheidungen
direkt bestimmt:

\begin{description}
  \item[Binärer Greifer] Über NeuraPy lässt sich der Greifer nur öffnen und schließen; eine
    Rückmeldung der Backenstellung gibt es nicht. Der Greiferzustand geht deshalb als
    kommandierter Wert $\{0, 1\}$ in Zustand und Aktion ein, und nach jedem Greiferbefehl
    wird eine feste Wartezeit von \SI{0.5}{\second} eingehalten, bevor sich der Arm weiter
    bewegt.
  \item[Schnittstelle zur Steuerung] NeuraPy spricht die Steuerung als Client über
    TCP/IP an. Jeder Aufruf ist ein Netzwerk-Roundtrip mit eigener Latenz
    (\cref{ch:ausfuehrung}).
  \item[Kinematik in der Steuerung] Inverse und Vorwärtskinematik stehen nur als Aufruf an
    die Steuerung zur Verfügung. Das mitgelieferte \ac{urdf}-Modell ließ sich nicht
    verbindlich mit der Steuerung abgleichen (\cref{sec:kinematik}).
  \item[Gemeinsam genutzte Anlage] Roboter und Draufsichtkamera werden auch von den Modulen
    1 und 2 benutzt. Die Draufsichtkamera ist an einer Stange montiert und kann zwischen
    Sitzungen verschoben werden; Prüfvorrichtung und Kamerastange liegen nah am
    Arbeitsbereich.
\end{description}

% ==============================================================================
\section{Festlegungen vor der ersten Aufnahme}
\label{sec:unumkehrbar}

Bei Behavior Cloning sind manche Entscheidungen nachträglich nicht korrigierbar, ohne alle
Daten neu aufzunehmen, weil die Policy die Bedingungen der Aufnahme mitlernt
(\cref{sec:weitere-huerden}). Diese Entscheidungen wurden deshalb vor der Aufnahme
ausdrücklich getroffen und begründet (\cref{tab:unumkehrbar}).

\begin{table}[htbp]
  \centering
  \caption[Festlegungen vor der ersten Aufnahme]{Festlegungen, die vor der ersten Aufnahme
  getroffen werden mussten, mit Begründung und Verweis auf den Abschnitt, der sie ausführt.}
  \label{tab:unumkehrbar}
  \small
  \begin{tabularx}{\textwidth}{@{}>{\raggedright\arraybackslash}p{30mm} >{\raggedright\arraybackslash}p{34mm} X@{}}
    \toprule
    Frage & Festlegung & Begründung \\
    \midrule
    Wie entstehen die Demonstrationen? & automatisch mit Bahnplaner und Rauschen &
      gleichmäßiges Tempo, keine Pausen durch Zögern, Korrekturbeispiele
      (\cref{sec:noise-umsetzung}) \\
    Aktionsraum & Gelenkwinkel, absolut, plus Greifer &
      entspricht 1:1 dem Befehl \code{servo_j}; absolute Sollwerte passen zu den
      asymmetrischen Datenpaaren (\cref{eq:asymm}) \\
    Zustandsraum & Gelenkwinkel, \acs{tcp}-Pose und Greifer (14 Werte) &
      die \acs{tcp}-Pose kostet kaum Speicher und wird für Auswertung und
      Kollisionsprüfung ohnehin gebraucht \\
    Takt & \SI{15}{\hertz} &
      im üblichen Bereich von LeRobot, mit Kameras und Steuerung sicher erreichbar \\
    Geschwindigkeit & fester Wert des Planers &
      die Policy lernt das Tempo aus dem Abstand zweier Aktionen mit \\
    Kameras & Handgelenk- und Draufsichtkamera &
      die Draufsicht liefert den Überblick, die Handgelenkkamera die Feinlage
      (\cref{sec:kameras}) \\
    Bildgröße & $240 \times 320$, RGB & gleich bei Aufnahme und Ausführung \\
    Rektifizierung & keine (revidiert) & siehe \cref{sec:kameras} \\
    \bottomrule
  \end{tabularx}
\end{table}

Um zu verhindern, dass ein Datensatz mit abweichenden Einstellungen unbemerkt verwendet wird,
trägt jeder Datensatz eine Schemaversion. Jede Änderung an Zustand, Aktion oder Bildgröße
erhöht sie, und Training und Ausführung brechen ab, wenn Modell und aktuelle Einstellungen
nicht zusammenpassen (\cref{sec:datensatz}).
